\documentclass[11pt]{article}
\usepackage{amsmath,amssymb,amsthm}
\newtheorem{theorem}{Theorem}
\title{Three complete streams can make an empirical risk cap hard to check}
\begin{document}
\maketitle

Two partially linked streams admit sharp upper and lower empirical
failure counts by weighted assignment. A third stream changes the
problem even when each non-anchor stream can be linked independently
to the anchor. The difficulty is not validating a candidate join.
It is choosing the joins whose failures cover the most or fewest
of the same episodes.

\paragraph{Interface.}
An anchor stream has one record per episode. Every other complete
stream has the same number \(n\) of records, one per episode, and
has a source-justified bipartite graph of permitted record-to-anchor
links. Each of its perfect matchings is a possible linkage, and
any collection of those matchings is jointly possible; there
are no additional cross-stream source constraints. Assume each
graph has a perfect matching. The audit asks for the union of
binary failures across streams. This is an empirical sample
calculation under a declared linkage model, not a time model or
a probability estimate. Arbitrary graphs need not be realizable
in the bounded-clock regime of the earlier recorder study.

\begin{theorem}
With one anchor and two other streams, deciding whether some
linkage attains at least a given empirical failure count is
NP-complete. Therefore deciding whether every linkage meets an
upper cap is coNP-complete, even at any fixed rational fraction
in \((0,1)\). Deciding whether some linkage attains at most
a given count is also NP-complete. These claims hold when the
anchor has no failures and each graph is independently feasible.
\end{theorem}
\begin{proof}
Matchings are polynomial witnesses for the existential decisions.
Reduce from perfect three-dimensional matching: sets \(X,Y,Z\)
each have size \(q\), and the \(n\) given triples are the anchor
episodes, initially all successful. Stream B has a failure
record for each \(x\in X\), allowed exactly at that element's
incident triples. For each \(y\in Y\) there are
\(\deg(y)-1\) successful records, confined to the triples
incident to \(y\). Thus a possible set \(S\) of failure episodes
in B has exactly one triple per \(x\) and per \(y\).
Conversely any such \(S\) extends to a valid B matching by pairing
the other records inside their \(y\)-classes.

For the lower-count question, stream C has one failure record
per \(z\), allowed exactly at its incident triples, and \(n-q\)
unconstrained successful records. A possible failure set \(T\)
has one triple per \(z\), and both \(|S|\) and \(|T|\) equal \(q\).
Hence \(|S\cup T|\le q\) iff \(S=T\), a perfect matching in
the supplied triples.

For the upper-count question, instead give the constrained C
records successful flags and the unconstrained records failure
flags. Its success set \(T\) again picks exactly one triple per
\(z\). Now the overall count is
\(|S\cup(\{1,\ldots,n\}\setminus T)|=n-|T\setminus S|\).
It is \(n\) iff \(S=T\). An upper cap of \(n-1\) therefore
holds iff the original triple instance has no perfect matching.
For a fixed rational cap fraction \(0<\alpha<1\) and \(n\ge2\),
add isolated matched safe episodes to bring the total to
\(\lceil(n-1)/\alpha\rceil\); the allowed integer count is
still \(n-1\).

If a triple class is empty, or B cannot be perfectly matched,
the original instance is negative; both can be checked before
construction. Map those cases to a fixed negative, individually
feasible instance in the appropriate reduction. For example,
use the four even-parity triples on \(\{0,1\}^3\); no pair is
complementary, B can match the two \(x\) and two \(y\) classes,
and C can match its \(z\) classes. This preserves the independent
feasibility promise.
\end{proof}

\paragraph{An exact easier upper bound under a different source law.}
If every successful non-anchor record can be assigned to any
episode, the maximum count is polynomial even for many streams.
Build one graph of all failure records from the other streams
and all episode positions not already failed at the anchor.
The answer is the anchor's count plus the maximum cardinality
matching in this graph. A joint labeling cannot do better:
choose one representative failure record for each additionally
failed position. Conversely, start with such a maximum matching
and restrict it to each stream. Augment that stream's partial
matching of failures until all its failure records are covered.
This is possible by feasibility, and augmenting paths keep every
previously assigned position occupied by a failure, even if the
specific record changes. Assign the unconstrained successes to
the remaining positions. Every position counted by the global
matching still fails. It is the constraints on otherwise healthy
records, not merely there being three data sources, that matter
to the hardness construction.

\paragraph{Evidence and scope.}
The replay checks all 256 hypergraphs on two elements per part,
including 193 with separately feasible constructed graphs. It
checks both extrema via exact feasible-set enumeration and checks
the easier upper-count algorithm on 380 separate seeded cases.
These are authored graph comparisons. They do not show that a
particular sensor source admits those graphs, or establish sensor
fault support, population risk, or full R4 acceptance. Direct
assignment still applies to the two-stream problem.
\end{document}
